\begin{center}
    {\LARGE \textbf{2018无机化学}} \\[2ex]
    {\large 时量180分钟 \quad 满分100分}
\end{center}

\vspace{0.5cm}
\noindent\textbf{注意事项:}
\begin{enumerate}[label=\arabic*., parsep=0pt, itemsep=0pt]
    \item 答题前,考生先将自己的姓名、学号、考场号、座位号填写在答题卡上,将条形码准确粘贴在考生信息条形码粘贴区。
    \item 考试结束后,将本试卷与答题纸一并收回。
    \item 回答各个问题时,将答案写在答题纸上,写在本试卷上无效。
\end{enumerate}

\vspace{0.5cm}
\noindent\textbf{一、选择题(20 pts.)}

\begin{enumerate}[label=\arabic*., itemsep=1.5ex]
    \item 下列硫化物中既能溶于$\ce{Na2S}$溶液又能溶于$\ce{Na2S2}$溶液的是
    \begin{tasks}(4)
        \task $\ce{ZnS}$
        \task $\ce{As2S3}$
        \task $\ce{HgS}$
        \task $\ce{CuS}$
    \end{tasks}
    \item 将$\ce{NCl3}$通入$\ce{NaOH}$溶液，得到产物是
    \begin{tasks}(4)
        \task $\ce{NH3}$和$\ce{NaCl}$
        \task $\ce{NH3}$和$\ce{NaClO}$
        \task $\ce{NO2}$和$\ce{NaCl}$
        \task $\ce{NH4Cl}$
    \end{tasks}
    \item 下列配离子中，构型会发生畸变的是
    \begin{tasks}(4)
        \task $\ce{[Cr(H2O)6]^3+}$
        \task $\ce{[Fe(CN)6]^4-}$
        \task $\ce{[Co(CN)6]^4-}$
        \task $\ce{[Ni(NH3)6]^2+}$
    \end{tasks}
    \item 下列化合物中，既溶于$\ce{HNO3}$溶液又溶于$\ce{NaOH}$溶液的是
    \begin{tasks}(4)
        \task $\ce{PbCrO4}$
        \task $\ce{BaCrO4}$
        \task $\ce{Ag2CrO4}$
        \task $\ce{BaSO4}$
    \end{tasks}
    \item 下列溶液中，加入$\ce{AgNO3}$溶液生成的沉淀颜色最浅的是
    \begin{tasks}(4)
        \task $\ce{H3PO3}$
        \task $\ce{NaH2PO4}$
        \task $\ce{Na2HPO4}$
        \task $\ce{NaPO3}$
    \end{tasks}
    \item 下列氢卤酸中，酸性最强的是
    \begin{tasks}(4)
        \task $\ce{HF}$
        \task $\ce{HCl}$
        \task $\ce{HBr}$
        \task $\ce{HI}$
    \end{tasks}
    \item 下列化合物中，碱性最弱的是
    \begin{tasks}(4)
        \task $\ce{NH3}$
        \task $\ce{N2H4}$
        \task $\ce{NH2OH}$
        \task $\ce{N3H}$
    \end{tasks}
    \item 下列物质中不能将$\ce{Mn^2+}$转化为$\ce{MnO4^-}$的是
    \begin{tasks}(4)
        \task $\ce{PbO2}$
        \task $\ce{H2O2}$
        \task $\ce{(NH4)2S2O8}$
        \task $\ce{H5IO6}$
    \end{tasks}
    \item 下列化合物中。都能与水发生氧化还原反应的是
    \begin{tasks}(4)
        \task $\ce{NH3}$,$\ce{SiH4}$
        \task $\ce{PH3}$,$\ce{B2H6}$
        \task $\ce{B2H6}$,$\ce{SiH4}$
        \task $\ce{PH3}$,$\ce{SiH4}$
    \end{tasks}
    \item 将标准电极中$\ce{H+}$浓度和$\ce{H2}$的分压减小为原数值的一半，其电极电势为
    \begin{tasks}(4)
        \task $0V$
        \task $0.009V$
        \task $-0.004V$
        \task $-0.009V$
    \end{tasks}
\end{enumerate}

\vspace{0.5cm}
\noindent\textbf{二、分析题(30 pts.)}

\begin{enumerate}[label=\arabic*., start=1, itemsep=1.5ex]
    \item $\ce{LiNO3}$受热分解生产气体为棕色，$\ce{NH4NO3}$受热分解生成的气体为无色。
    \item 将$\ce{KMnO4}$晶体的热分解产物放入试管中，加少许水时现象为试管上部为绿色溶液，但加入大量水后试管上部为紫色溶液。
    \item 向盛有$\ce{Na2S2O3}$溶液的试管中加一滴$\ce{AgNO3}$溶液，生成白色沉淀，摇动试管则沉淀立即溶解;向盛有$\ce{AgNO3}$溶液的试管中加一滴$\ce{Na2S2O3}$溶液，生成白色沉淀，摇动试管沉淀不溶解，之后沉淀逐渐变黄，变棕，变黑。
    \item 用浓$\ce{HCl}$处理$\ce{Co2O3}$得到蓝色溶液，加水稀释后溶液变为粉红色。
    \item 向酸性$\ce{KIO3}$的淀粉溶液中滴加$\ce{Na2SO3}$溶液，溶液先变蓝后变无色;向$\ce{Na2SO3}$的淀粉溶液中滴加酸性$\ce{KIO3}$溶液，溶液先无色后变蓝。
    \item 向酸性$\ce{H2O2}$溶液中滴加$\ce{K2Cr2O7}$溶液得到蓝色溶液，向碱性$\ce{H2O2}$溶液中滴加$\ce{K2Cr2O7}$溶液得到黄色溶液。
\end{enumerate}

\vspace{0.5cm}
\noindent\textbf{三、简答题(54 pts.)}

\begin{enumerate}[label=\arabic*., start=1, itemsep=1.5ex]
    \item 说明下列各实验中浓$\ce{HCl}$的作用
    \begin{enumerate}[label=(\arabic*), itemsep=1ex]
        \item 配置$\ce{SnCl2}$溶液时，将$\ce{SnCl2(s)}$溶于浓$\ce{HCl}$中而后加水稀释。
        \item 加热条件下用$\ce{MnO2}$与浓$\ce{HCl}$溶液制取氯气。
        \item 用浓$\ce{HCl}$溶液配置的王水溶解金。
    \end{enumerate}
    \item 比较下列各对配合物稳定性，并简要说明原因
    \begin{enumerate}[label=(\arabic*), itemsep=1ex]
        \item $\ce{[Co(NH3)6]^2+}$与$\ce{[Co(NH3)6]^3+}$
        \item $\ce{[Cu(NH3)4]^2+}$与$\ce{[Cu(en)2]^2+}$
        \item $\ce{[HgI4]^2-}$与$\ce{[HgCl4]^2-}$
        \item $\ce{[Cu(NH3)4]^2+}$与$\ce{[Zn(NH3)4]^2+}$
    \end{enumerate}
    \item 给出下列分子或离子的几何构型和中心原子杂化类型
    \begin{enumerate}[label=(\arabic*), itemsep=1ex]
        \item $\ce{NO2}$
        \item $\ce{SOCl2}$
        \item $\ce{ClF3}$
        \item $\ce{[Ni(CN)4]^2-}$
        \item $\ce{XeOF4}$
        \item $\ce{SF4}$
    \end{enumerate}
    \item 比较大小（用>或<号表示）。
    \begin{enumerate}[label=(\arabic*), itemsep=1ex]
        \item 熔点: $\ce{NaF}$,$\ce{NaCl}$,$\ce{NH3}$,$\ce{PH3}$
        \item 分子极性: $\ce{NH3}$,$\ce{NF3}$,$\ce{O2}$,$\ce{O3}$
        \item 热稳定性$\ce{PbCl4}$,$\ce{PbCl2}$,$\ce{MnO2}$,$\ce{PbO2}$
        \item 溶解度: $\ce{LiF}$,$\ce{NaF}$,$\ce{KClO4}$,$\ce{NaClO4}$
        \item 键级: $\ce{O2^-}$,$\ce{O2^{2-}}$;$\ce{O2}$,$\ce{O2^+}$
    \end{enumerate}
    \item 写出具有下列电子构型的元素名称，符号及在周期表中位置
    \begin{enumerate}[label=(\arabic*), itemsep=1ex]
        \item $3d^54s^1$
        \item $4d^45s^1$
        \item $4f^75d^16s^2$
        \item $4d^{10}$
        \item $5s^25p^1$
    \end{enumerate}
    \item 已知配离子$\ce{[Co(NH3)6]^3+}$的$K_{\text{稳}}=1.58\times10^{35}$，电对$\ce{(Co^{3+}/Co^{2+})}$的$E^\ominus=1.92V$，电子对$\ce{[Co(NH3)6]^3+/Co(NH3)6^2+}$的$E^\ominus=0.108V$，试求配离子$\ce{[Co(NH3)6]^{2+}}$的$K_{\text{稳}}$
    \item 用姜-泰勒效应讨论$\ce{[PtCl4]^2-}$的结构和磁性。
\end{enumerate}

\vspace{0.5cm}
\noindent\textbf{四、制备与分离(20 pts.)}

\begin{enumerate}[label=\arabic*., start=1, itemsep=1.5ex]
    \item 用反应方程式表示下列制备过程
    \begin{enumerate}[label=(\arabic*), itemsep=1ex]
        \item 以碳酸钠和硫磺为原料制备硫代硫酸钠
        \item 由$\ce{Cl2}$制备$\ce{O2}$
    \end{enumerate}
    \item 设计方案分离$\ce{Zn^2+}$,$\ce{Cr^3+}$,$\ce{Mn^2+}$,$\ce{Al^3+}$
    \item 以$\ce{Fe}$屑为原料制备赤血盐。
\end{enumerate}

\vspace{0.5cm}
\noindent\textbf{五、鉴别与推断(26 pts.)}

\begin{enumerate}[label=\arabic*., start=1, itemsep=1.5ex]
    \item 设计实验以验证铅丹中铅不同氧化态，给出反应方程式
    \item 白色固体$\ce{A}$加入水中生成白色沉淀$\ce{B}$，取上层澄清溶液加入$\ce{AgNO3}$有白色沉淀$\ce{C}$生成，$\ce{C}$不溶于$\ce{HNO3}$和$\ce{NaOH}$溶液，而易溶于氨水，$\ce{B}$不溶于$\ce{NaOH}$溶液，溶于强酸，$\ce{A}$溶于盐酸溶液后加入$\ce{NaOH}$溶液生成白色沉淀$\ce{D}$，再加入$\ce{NaClO}$溶液有土黄色沉淀$\ce{E}$生成，酸性条件下，将$\ce{E}$与$\ce{MnSO4}$混合，溶液变成紫红色，$\ce{A}$由盐酸溶液与$\ce{H2S}$溶液反应生成黑色沉淀$\ce{F}$ ，向$\ce{SnCl2}$溶于$\ce{NaOH}$溶液中投入少量$\ce{A}$，生成黑色沉淀$\ce{G}$，确定$\ce{ABCDEFG}$所代表的化学式，写出$\ce{D\rightarrow E}$的羟化反应的相关反应方程式
    \item 粉红色晶体$\ce{A}$受热脱水失重约45.5\%得到白色粉末$\ce{B}$，$\ce{B}$溶于水得到无色溶液，再加入$\ce{BaCl2}$有不溶于硝酸的白色沉淀$\ce{C}$，$\ce{A}$ 溶于水后加入$\ce{NaOH}$溶液，生成白色沉淀$\ce{D}$，$\ce{D}$在空气中缓慢变成棕黑色，说明有$\ce{E}$ 产生，$\ce{E}$与稀硫酸和双氧水混合液作用溶解生成无色溶液$\ce{F}$，$\ce{E}$与$\ce{KOH}$与$\ce{KClO3}$共熔，生成绿色产物$\ce{G}$，将$\ce{G}$投入稀硫酸有紫红色溶液$\ce{H}$生成，请给出$\ce{A}$、$\ce{D}$、$\ce{E}$、$\ce{G}$、$\ce{H}$的化学式及$\ce{E\rightarrow G}$转化反应的方程式
    \item 将化合物$\ce{A}$溶于水后加入$\ce{NaOH}$溶液有黄色沉淀$\ce{B}$生成，$\ce{B}$不溶于过量的$\ce{NaOH}$溶液，$\ce{B}$溶于$\ce{HCl}$溶液得无色溶液，向该溶液中滴加少量$\ce{SnCl2}$溶液有白色沉淀$\ce{C}$ 生成，向$\ce{A}$ 的水溶液中滴加$\ce{KI}$溶液得红色沉淀$\ce{D}$，$\ce{D}$可溶于过量$\ce{KI}$溶液得无色溶液，向$\ce{A}$的水溶液中加入$\ce{AgNO3}$溶液有白色沉淀$\ce{E}$生成，$\ce{E}$不溶于$\ce{HNO3}$溶液但可溶于氨水生成$\ce{F}$，写出$\ce{A-F}$代表的物质的化学式，并写出生成$\ce{C}$的化学方程式。
\end{enumerate}